\chapter{List of exercises}

\begin{dang}
\textbf{Exercise 1}. Let $U=k^2-\{0\}$, show that the set of regular functions $U\to k$ are polynomials, then 
conclude that $U$ is not an affine variety.    
\end{dang}
\textbf{Solution}. Let $\varphi\in\mathcal{O}_U(U)$ be a regular function, then 
$\varphi|_{D_{k^2}(X)}=f/X^m$ and $\varphi|_{D_{k^2}(Y)}=g/Y^n$ and on the intersection $D_{k^2}(X)\cap D_{k^2}(Y)=D_{k^2}(XY)$, we have 
\begin{equation}
    Y^nf - X^mg = 0\quad(*).
\end{equation}
Since there exist an infinity number of points $p\in D_{k^2}(XY)$ (assuming $k$ is algebraically closed), 
we also have the equality $(*)$ in $k[X,Y]$. Now, let $f=\sum a_{ij}X^iY^j$ and 
$g=\sum b_{ij} X^iY^j$. Then, equality $(*)$ means 
\begin{equation}
    a_{i, j-n} = a_{i-m,j},\;\forall i\geq m, j\geq n.
\end{equation}  
Therefore $f/X^m = g/Y^n\in k[X,Y]$ and $\varphi$ is a polynomial, because $U=D_{k^2}(X)\cup D_{k^2}(Y)$. In conclusion,
we proved that the inclusion $i:U\hookrightarrow k^2$ induces a ring isomorphism
\begin{equation}
    i^\#(k^2): k[X,Y]\to \mathcal{O}_U(U):f\mapsto f|_U\; (**).
\end{equation}
Now assume $U=V(S),\;S\subset k[X,Y]$, then $f\in S\Rightarrow f|_U=0$ and $f=0$ because of $(**)$. So 
$S=\{0\}$ and $U=k^2$, which is a contradiction.\qed



